\section{準備}\label{sec:prelim}

\subsection{変数}
\begin{center}
\small
\begin{tabular}{>{$}l<{$}p{9cm}}
\toprule
\text{記号} & 意味 \\
\midrule
q,\ t & refined CS の変数。$r = q^{1/2}$, $s = t^{1/2}$ \\
U = t^N & SO$(2N)$ の Kähler 変数（$N$ は D$_N$ のランク）。大きな $N$ では形式変数 \\
\nu = U t^{-1/2} & Mariño の $\nu$（$t = q$ で $\nu = q^{N-1/2}$） \\
\tnu = \nu\, (q/t)^{1/2} & ずれた Kähler 変数 \\
A = \nu^2 = t^{N_A},\ N_A = 2N - 1 & 被覆 U$(N_A)$ の Kähler 変数 \\
(\bq, \bt, \ba) = (t^{1/2}, -(q/t)^{1/2}, \nu) & 太字の変数（DGR \cite{DGR05}・\cite{AgSc12} と同じ）。$q^{1/2} = -\bt\bq$ \\
V = (qt)^{1/4},\ W = (q/t)^{1/4} & $q^{1/2} = VW$, $t^{1/2} = V/W$ \\
z_q = q^{1/2} - q^{-1/2},\ z_t = t^{1/2} - t^{-1/2} & $z_q^+ = q^{1/2} + q^{-1/2}$ \\
G = z_q z_t & $G$ 展開の変数（KN \cite{KN17}） \\
\sigma\colon q \mapsto t^{-1},\ t \mapsto q^{-1} & $\nu$ を固定。$G$ は不変 \\
\bottomrule
\end{tabular}
\end{center}
unrefined の Mariño の変数は $t_M = q^{1/2}$、$z = z_q$ である。

\subsection{refined SO$(2N)$ の不変量}
AgSc \cite{AgSc12} の SO$(2N)$ の refined Chern--Simons 理論では、表現 $\lambda$ の Wilson ループは D$_N$ 型の Macdonald 多項式で与えられる。
% [S1]
\begin{compresult}
AgSc の refined unknot $\dim_\lambda$ は、Koornwinder 多項式 $K_\lambda(x; q, t; 1, -1, q^{1/2}, -q^{1/2})$ の主特殊化 $x = (1, t, \dots, t^{N-1})$ に等しい（6 つの $\lambda$。van Diejen--Sahi の評価式と、$q$-差分作用素で作った多項式（$N = 2, 3$）で確認）。
$\dim_\lambda$ は $U$ の Laurent 多項式ではなく、例えば [2] は $(1 - q t^{N-1})$、[1,1] は $(1 - q t^{2N-3})$ を分母にもつ（$t = q$ では約分される）。
\end{compresult}

% [S2]
\begin{remark}
Mariño 型の構成（bosonic な Cauchy 核、$\tfrac12 f^{\cov}$、$\sum_{d\ \mathrm{odd}}$）をそのまま refined にすると、unknot の 2 箱で $U$ の極が打ち消さない。
composite $([1],[1])$ の寄与は $p_1^2$ 方向だけにあるが、D$_N$ の [2] と [1,1] の極の位置は違う（$U = t/q$ と $U^2 = t^3/q$）。
重み 5 通りと基底 3 通りのどの組でもこれは解消しない。
\end{remark}

\subsection{双対 Cauchy 核と生成関数}
% [S3]
Mimachi の双対 Cauchy 恒等式（Koornwinder 多項式）
\[
\sum_{\lambda \subseteq m^n} (-1)^{|\lambda|} K_{m^n - \lambda}(x; q, t)\, K_{\lambda'}(y; t, q) = \prod_{i, j} x_i^{-1}(1 - x_i y_j)(1 - x_i y_j^{-1})
\]
では、ホロノミー側に AgSc の D$_n$ 型（$\pm1, \pm q^{1/2}$）、probe 側に $(q', t') = (t, q)$ に対する $(\pm 1, \pm t'^{1/2})$ 型が現れる。
この恒等式は $(n, m) = (1,1), (2,1), (1,2), (2,2), (2,3), (3,2)$ の 6 通りで確かめた。
これを refined orientifold の Ooguri--Vafa 核と読み、probe 変数 $Y = (y_1, y_2, \dots)$ に対して
\begin{equation}\label{eq:Z}
\tilde Z(Y) = \Bigl\langle \prod_{\alpha, i} (1 - y_\alpha x_i)(1 - y_\alpha x_i^{-1}) \Bigr\rangle_{\SO(2N)},\qquad
\log \tilde Z = \sum_{d \ge 1} \frac1d \sum_R f_R(q^d, t^d, U^d)\, s_R(Y^d)
\end{equation}
と置く。被覆 U$(N_A)$ についても、同じ核で $f^{\cov}_R$ を定める。
核の $|Y| \le 2$ の部分は、U$(N_A)$ の普通の表現と随伴 composite $(1;1)$ で書ける。3 次では composite $(2;1)$, $(1,1;1)$ も現れる。

\begin{definition}[クロスキャップ部分]
全 $d$ の定義で $g_R = f_R - \tfrac12 f^{\cov}_R$ と置く。
hatted 不変量は $\hat g = M(q^{-1})^{-1} g$ である。ここで $M_{RS} = \sum_\mu \chi_R(\mu)\chi_S(\mu) z_\mu^{-1}\prod_i (t_M^{\mu_i} - t_M^{-\mu_i})/(t_M - t_M^{-1})$（$t_M = q^{1/2}$）は Mariño の $M$ で、KN の $M(x)$ の $x = q^{-1}$ に等しい。
\end{definition}

% [S8]
\begin{remark}
reformulation に $M(t)$ や $M(t^{-1})$ を使うと、Hopf 絡み目で Laurent にならない。$M(q)$ は $R \leftrightarrow R^T$ を入れ替える。
KN の二つの枠（q 枠：Macdonald の Cauchy 核と $M(t)$、t 枠：双対 Cauchy 核と $M(q^{-1})$）の hatted 不変量は、2 箱で一致する（§\ref{sec:generalN}）。
\end{remark}

\subsection{一様な規格化}\label{sec:uniform}
入力の U$(n)$ の refined 不変量（KN の (N1)+(N2) で規格化したもの）は、行列式表現と荷電共役について $(q/t)$ の冪だけずれる。
そこで、[1] の superpolynomial の $\delta$ 次数 $\delta = \ba + \bq/2 - \bt$ を使い、
\begin{equation}\label{eq:uniform}
W_\mu = (t/q)^{\delta |\mu|/2}\, R_\mu\, P_\mu(t^\rho)
\end{equation}
と置く（$T(2,2m+1)$ では $\delta = m$）。

% [S5][S6]
\begin{compresult}\label{cr:uniform}
\begin{enumerate}
\item 規格化 \eqref{eq:uniform} では、行列式のずらし $\mu \to \mu + 1^n$ と荷電共役で $W$ が正確に不変になる。検査した範囲は trefoil の 16 組（U(2)〜U(5)）、$T(2,5)$ と $T(2,7)$ の各 13 組、二重ツイスト結び目 7 個である。
\item 固定 $N$（$N = 2, 3$、$|R| \le 3$）で framing の単項式を「$N$ に依存する極の打ち消し」と「$t = q$ で 1」で決めると、SO 側・被覆側ともすべて 1 になる（$T(2,3), T(2,5), T(2,7)$）。この走査では $(t/q)^{\beta/2}$ の族の中で $\beta$ が一意に決まる。
\item DAHA-Jones 多項式（鉛直 framing）とは、framing 固有値の冪を除いて一致する（§\ref{sec:generalN}）。
\item $t = q$ では Mariño の標準 framing そのものになる。
\end{enumerate}
\end{compresult}

% [S7]
\begin{hypothesis}
一様な規格化 \eqref{eq:uniform} は framing 0 の値である。
\end{hypothesis}
以後、SO 側・被覆側ともこの規格化を使う。
